\documentclass{article}
\usepackage[T1]{fontenc}
\usepackage{lmodern}
\usepackage{graphicx}
\usepackage{amsmath,amssymb}
\usepackage[a4paper,margin=2cm]{geometry}
\usepackage[absolute,overlay]{textpos}
\setlength{\TPHorizModule}{1mm}
\setlength{\TPVertModule}{1mm}
\usepackage[indonesian]{babel}
\usepackage{hyperref}
\setlength{\emergencystretch}{2em}
% unit: Halaman 99

\begin{document}

% === Halaman 99 (python/native) ===

• Cara menghitung matriks invers 2 x 2: 

   K =                          K-1 = 

                                        =

  Contoh: K =               


   det(K) = (3)(9) – (15)(10) = 27 – 150  = –123 mod 26 = 7                    

99













d

c

b

a













−

−

a

c

b

d

K)

det(

1













−

−

−

a

c

b

d

bc

ad

1













9

15

10

3

\end{document}
